\documentclass[10pt,twocolumn]{article}
\usepackage[T1]{fontenc}
\usepackage[utf8]{inputenc}
\usepackage{lmodern}
\usepackage[margin=1.8cm,top=2.4cm,bottom=2cm,columnsep=0.6cm]{geometry}
\usepackage{fancyhdr}
\usepackage{titlesec}
\usepackage{booktabs}
\usepackage{amsmath,amssymb,amsfonts}
\usepackage{microtype}
\usepackage{xcolor}
\usepackage{mdframed}
\usepackage{parskip}
\usepackage{enumitem}
\usepackage{hyperref}

\definecolor{rulered}{RGB}{180,30,30}
\definecolor{boxgray}{RGB}{245,245,245}
\definecolor{keywordgray}{RGB}{100,100,100}

\pagestyle{fancy}
\fancyhf{}
\fancyhead[L]{\textsc{\small Claude Mag}}
\fancyhead[C]{\small\textit{Machine Learning Research Digest}}
\fancyhead[R]{\small 2026-11-02}
\fancyfoot[C]{\small\thepage}
\renewcommand{\headrulewidth}{0.8pt}

\titleformat{\section}{\normalsize\bfseries\color{rulered}}{}{0pt}{}%
  [\vspace{-4pt}\color{rulered}\rule{\columnwidth}{0.4pt}]
\titlespacing{\section}{0pt}{8pt}{4pt}

\hypersetup{colorlinks=true,urlcolor=rulered,linkcolor=black}

\begin{document}
\setlength{\parindent}{0pt}
\setlength{\parskip}{4pt}

{\small\color{keywordgray}\textbf{Keywords:}
  flow matching \textbullet{}
  language generation \textbullet{}
  diffusion language models \textbullet{}
  fast inference}\par\vspace{4pt}

{\LARGE\bfseries\raggedright
  Rectified Language Flow}\par\vspace{4pt}

{\large\itshape\raggedright
  Straight-Path Flow Matching in Token Embedding Space Beats
  Autoregressive Baselines at 15--19$\times$ Inference Speedup}\par\vspace{6pt}

{\small\textbf{By} Baiyu Su, Runlong Liao, Lizhang Chen, Qiang Liu}\par\vspace{2pt}

{\small\color{keywordgray}arXiv:2610.33811 \textbullet{}
  cs.AI / cs.LG \textbullet{} University of Texas at Austin}
\par\vspace{2pt}\noindent{\color{rulered}\rule{\columnwidth}{0.6pt}}\vspace{6pt}

Rectified Language Flow (ReLF) generates text by training a flow model
to move noise continuously toward token embeddings along straight
(rectified) paths, enabling accurate generation in as few as one to
four function evaluations. At batch size 1 it produces 1,024-token
sequences 15--19$\times$ faster than a matched autoregressive baseline,
while achieving lower generative perplexity on OpenWebText and
outperforming all prior continuous-flow language models.

\begin{mdframed}[backgroundcolor=boxgray,linecolor=rulered,linewidth=1pt,
    innerleftmargin=6pt,innerrightmargin=6pt,
    innertopmargin=6pt,innerbottommargin=6pt]
\textbf{\small AT A GLANCE}\par\vspace{4pt}
\begin{description}[leftmargin=0pt,labelsep=4pt,itemsep=2pt,parsep=0pt,
    font=\small\bfseries]
  \item[Problem:] Autoregressive models are slow at inference due to
    serial token generation; discrete diffusion models require many
    costly denoising steps over a large discrete space.
  \item[Why it matters:] Fast, high-quality text generation is
    critical for interactive AI, serving efficiency, and real-time
    applications such as code completion and dialogue.
  \item[Method:] Continuous flow matching in Hellinger-SVD token
    embeddings with self-conditioning guidance, regularized reflow for
    distillation, and singular logits for discrete commitment.
  \item[Result:] Lower perplexity than autoregressive baselines on
    OpenWebText; 15--19$\times$ faster generation at batch size 1
    with 4 function evaluations.
  \item[Evidence:] OpenWebText language modeling benchmarks across
    three model sizes; ablation over number of function evaluations.
\end{description}
\end{mdframed}

\section{The Big Picture}

Large language models have converged on autoregressive (AR) generation:
at each step the model conditions on all previous tokens and samples
the next one. This left-to-right serial process is the primary
bottleneck at inference time---each token requires a full forward pass,
and a sequence of $T$ tokens requires $T$ passes.

Diffusion language models offer an alternative: start from a noisy
sequence and iteratively refine it over $K$ denoising steps. Because
denoising steps can be parallelized across positions, the number of
forward passes is $K \ll T$ for long sequences. But discrete diffusion
models operate over a large categorical space at each position,
introducing its own inefficiencies, and $K$ is typically still in the
tens to hundreds.

ReLF takes a third approach: flow matching in a continuous embedding
space. A flow model is trained to transport a simple noise distribution
$p_0$ (e.g., Gaussian) to the distribution of token embeddings $p_1$.
The rectified flow objective enforces straight paths between noise and
data, enabling accurate generation with as few as 1--4 function
evaluations.

\section{Why Existing Approaches Fall Short}

\textbf{Autoregressive models} require $T$ serial forward passes for a
length-$T$ sequence. Speculative decoding and other parallelization
heuristics reduce this in practice but don't change the fundamental
asymptotic cost.

\textbf{Discrete diffusion models} work directly in token space,
avoiding the need to commit to embeddings. However, the discrete state
space is large (vocabulary $\sim 50\text{k}$), and the denoising
objective must define meaningful noise over discrete categories. Prior
models report 32--64 steps to achieve quality competitive with AR
models.

\textbf{Prior continuous-flow language models} (ELF, DiLaDiff) use
learned continuous embeddings that are unstable to train and require
reconstruction losses that decouple the flow from the discrete output
space. Committing a continuous embedding to a discrete token at
generation time via nearest-neighbor lookup introduces discretization
error that accumulates over the sequence.

\section{Core Mechanism}

ReLF's four key components work together to address these issues:

\textbf{1. Hellinger-SVD Embeddings.} Token embeddings are initialized
from the SVD of the Hellinger-transformed co-occurrence matrix of the
training corpus. This produces a geometry in which frequent co-occurring
tokens are nearby, providing a meaningful flow target that is fixed
throughout training. Crucially, no embedding training is needed.

\textbf{2. Self-Conditioning Guidance.} At each training step the model
is first run with the corrupted input $x_t$ to produce a preliminary
estimate $\hat{x}_0$. On the same forward pass (with gradients detached),
$\hat{x}_0$ is concatenated to $x_t$ and passed to the model again.
A cross-entropy loss between $\hat{x}_0$ and the true $x_0$ trains the
guidance branch. This conditioning helps the model resolve ambiguous
intermediate states near $t = 0.5$.

\textbf{3. Regularized Reflow.} After initial training, a reflow stage
distills the learned curved flow into a straighter one by training on
$(x_0, x_1)$ pairs sampled from the original model's trajectories.
A regularization term penalizes degenerate straight lines that pass
outside the data manifold, preventing the reflow from creating diversity
collapse.

\textbf{4. Singular Logits.} As $t \to 1$ the flow approaches the data
manifold and must commit to a discrete token. Singular logits is a
differentiable mechanism that computes token logits from the flow's
velocity direction at $t = 1$, rather than from the embedding value
itself, enabling sharp, gradient-compatible token selection.

\section{Technical Formulation}

Let $x_0 \sim \mathcal{N}(0, I)$ be noise and $x_1 \sim p_{\text{data}}$
be a token embedding sequence. The rectified flow objective minimizes:
\[
\mathcal{L}_{\text{flow}} = \mathbb{E}_{t, x_0, x_1}\!\left[
  \left\| v_\theta(x_t, t) - (x_1 - x_0) \right\|^2
\right]
\]
where $x_t = (1-t)x_0 + t\, x_1$ and $v_\theta$ is the learned velocity
field.

The self-conditioning guidance loss augments training with:
\[
\mathcal{L}_{\text{sc}} = -\log p_\theta(x_1 \mid \hat{x}_0, x_t, t)
\]
where $\hat{x}_0 = x_t + (1-t)\,v_\theta(x_t, t)$ is the preliminary
clean estimate.

The regularized reflow loss is:
\[
\mathcal{L}_{\text{reflow}} = \mathcal{L}_{\text{flow}}^{\text{straight}}
  + \lambda \cdot \mathbb{E}\!\left[\text{dist}(x_t^{\text{reflow}},
  \mathcal{M})\right]
\]
where $\mathcal{M}$ is an approximation of the data manifold and $\lambda$
is a regularization weight. Singular logits are defined as:
\[
\ell_v = W_{\text{vocab}}\, v_\theta(x_1, 1) \cdot \|W_{\text{vocab}}\|^{-1}
\]
mapping the terminal velocity to logits via the vocabulary matrix.

\section{Learning or Inference Procedure}

\textbf{Training:}
\begin{enumerate}[itemsep=2pt,parsep=0pt]
  \item Pre-compute Hellinger-SVD embeddings from corpus statistics.
  \item Train $v_\theta$ with $\mathcal{L}_{\text{flow}} + \alpha \cdot
    \mathcal{L}_{\text{sc}}$ for $N$ steps.
  \item Collect $(x_0, x_1)$ pairs from teacher model trajectories.
  \item Fine-tune with $\mathcal{L}_{\text{reflow}}$ for $M \ll N$
    additional steps.
\end{enumerate}

\textbf{Inference:} Sample $x_0 \sim \mathcal{N}(0, I)$, integrate
the velocity field over $K \in \{1, 2, 4, 8, 16, 64\}$ Euler steps,
apply singular logits at $t = 1$ to decode token distributions, and
sample.

\section{What the Guarantee Says}

The rectified flow objective guarantees that the expected trajectory
length $\mathbb{E}[\int_0^1 \|v_\theta\| \,dt]$ is minimized when
trajectories are straight, which corresponds to the lower bound on
the number of function evaluations required for a given discretization
error. While this is an asymptotic argument, the regularized reflow
empirically achieves near-straight paths and the 1-step sample quality
confirms the practical benefit. No PAC-style generalization bound is
provided; the guarantee is about transport efficiency, not sample quality.

\section{Experimental Findings}

\begin{itemize}[itemsep=2pt,parsep=0pt]
  \item \textbf{Perplexity:} On OpenWebText, with 64 function
    evaluations, ReLF achieves \textbf{lower perplexity than AR
    baselines} at matched parameter counts for all three model sizes
    (90M, 350M, 1.2B).
  \item \textbf{Speed:} At 4 function evaluations, ReLF generates
    1,024-token sequences \textbf{15--19$\times$ faster} than the AR
    baseline at batch size 1, with less than 0.5 perplexity units of
    degradation relative to the 64-step setting.
  \item \textbf{Comparison to ELF:} ReLF outperforms ELF at all model
    sizes by 1.2--2.8 perplexity points.
  \item \textbf{Diversity:} At matched perplexity, ReLF self-BLEU
    scores are lower (more diverse) than the AR baseline, suggesting
    the flow's sampling process does not collapse to near-deterministic
    outputs.
\end{itemize}

\section{Ablations and Interpretation}

\textbf{Hellinger-SVD vs.\ learned embeddings:}
Learned embeddings reduce final perplexity by 0.3 points on average
but add training instability; the fixed Hellinger-SVD embeddings are
more robust and are used in all reported experiments.

\textbf{Without self-conditioning:}
Removing self-conditioning increases perplexity by 1.8--3.4 points,
with the largest degradation at small $K$ (few steps). This confirms
that guidance is most critical in the low-step regime.

\textbf{Without reflow:}
Without the reflow stage, 4-step generation degrades by 5.1 perplexity
points relative to the 64-step setting. With reflow, the gap closes to
$<0.5$ points.

\textbf{Singular logits vs.\ nearest neighbor:}
Replacing singular logits with nearest-neighbor lookup in embedding
space increases perplexity by 0.7--1.1 points and degrades diversity.
The directional signal at $t=1$ is more discriminative than the
position signal.

\section*{Reference}

Baiyu Su, Runlong Liao, Lizhang Chen, Qiang Liu.
\textbf{Rectified Language Flow.}
arXiv:2610.33811. 2026.
\url{https://arxiv.org/abs/2610.33811}

\end{document}
